\section{Restriction and infinitesimal rigidity}
\label{sec:sr-restriction-rigidity}
We now prove the restriction and local tangent statements needed when
the integer comparisons are applied to their actual automorphic points.
The existence and automorphy of those points are separate inputs.


\subsection{Restriction to the infinite cyclotomic kernel}
\label{sec:sr-cyclotomic-restriction}
The ordinary type-I argument in
Section~\ref{sec:spectrum-open-irreducibility} proves irreducibility
on every open subgroup of the relevant absolute Galois group.
The cyclotomic kernel is closed and usually not open.
We give the additional group-theoretic argument needed to pass to it.

\begin{srlemma}\label{lem:sr-procyclic-restriction}
Let \(G\) be a profinite group, \(H\) a closed normal subgroup
with procyclic quotient, and \(L\) an algebraically closed
Hausdorff topological field of characteristic zero.
Let \(V\) be a finite-dimensional continuous \(L\)-representation
of \(G\).
Suppose \(V\) is irreducible on every open subgroup of \(G\).
Then \(V|_H\) is irreducible.
Finite-dimensional linear subspaces are understood with their
usual coordinate topology, so they are closed.
\end{srlemma}
\begin{proof}
Choose a nonzero \(H\)-submodule of smallest dimension.
It is simple. Every \(G\)-translate is again a simple
\(H\)-submodule, since \(H\) is normal.
Their sum is nonzero and \(G\)-invariant, hence is \(V\).
A finite subcollection already spans \(V\).
A sum of simple submodules is semisimple: add them successively,
omitting a summand if it is contained in the sum already chosen;
otherwise its intersection with that sum is zero by simplicity.
Thus \(V|_H\) is semisimple.

Group its simple summands into their isotypic components
\(V_1,\ldots,V_j\). These are intrinsic \(H\)-submodules and \(G\)
permutes them.
The stabilizer of \(V_i\) is closed: in a basis adapted to
\(V_i\), invariance is the vanishing of a fixed matrix block.
It has finite index, since there are finitely many components.
A closed finite-index subgroup of a profinite group is open,
as its complement is a finite union of closed cosets.
If there were more than one component, this open stabilizer
would preserve the proper nonzero subspace \(V_i\), contradicting
the hypothesis. Hence there is only one isotypic component.

Choose an \(H\)-module identification \(V=W\otimes_L M\), with
\(W\) simple and \(H\) acting trivially on \(M\).
Schur's lemma here is elementary: an \(H\)-equivariant endomorphism
of \(W\) has an eigenvalue over \(L\); the kernel after subtracting
that scalar is a nonzero \(H\)-submodule, so the endomorphism is scalar.
It follows, by writing a matrix on the multiplicity factors, that
\[
                 \operatorname{End}_H(V)
                           =1_W\otimes\operatorname{End}_L(M).
\]

Choose \(g\in G\) lifting a topological generator of \(G/H\).
Conjugation by \(g\) preserves the simple \(H\)-type \(W\).
Choose an intertwining isomorphism \(T\) on \(W\) for that conjugation.
After multiplying the action of \(g\) by \(T^{-1}\) on the first
factor, the resulting operator commutes with \(H\).
Thus the action of \(g\) is \(T\otimes B\) for some
\(B\in\mathrm{GL}(M)\).
Over the algebraically closed field \(L\), \(B\) has an eigenline
\(M_1\). The subspace \(W\otimes M_1\) is invariant under \(H\)
and \(g\), in both directions.

The subgroup generated by \(H\) and \(g\) is dense in \(G\).
Indeed its closure is compact, its image in \(G/H\) is closed and
contains a topological generator, and it contains \(H\);
its image is therefore all of \(G/H\), so the closure is all of \(G\).
The stabilizer of \(W\otimes M_1\) is closed by the same block-matrix
argument. It consequently equals \(G\).
Irreducibility of \(V\) now implies \(W\otimes M_1=V\), so
\(\dim M=1\). This proves irreducibility on \(H\).
\end{proof}

\begin{srproposition}\label{prop:sr-cyclotomic-restriction}
Let \(F\) be a number field and let \(\ell\) be odd.
A continuous representation over \(\overline{\mathbb Q}_\ell\)
which is irreducible on every open subgroup of \(G_F\)
is irreducible on \(G_{F(\mu_{\ell^\infty})}\).
This conclusion remains true after a finite extension of \(F\)
and after any scalar character twist.
\end{srproposition}
\begin{proof}
The quotient by the indicated kernel is the closed cyclotomic
image in \(\mathbb Z_\ell^\times\).
For odd \(\ell\),
\(\mathbb Z_\ell^\times\simeq\mu_{\ell-1}\times\mathbb Z_\ell\)
is procyclic; one element generates every finite quotient by
the Chinese remainder theorem, since the finite factor has
order prime to \(\ell\).
Every closed subgroup of a procyclic profinite group is procyclic.
Apply Lemma~\ref{lem:sr-procyclic-restriction}.
Every open subgroup of \(G_{F'}\), for \(F'/F\) finite, is also
open in \(G_F\), so the same hypothesis holds after this base change.
A scalar twist has precisely the same invariant subspaces as the
original representation on every subgroup.
\end{proof}

For the normalized ordinary type-I points, the open-subgroup
hypothesis is proved in
Section~\ref{sec:spectrum-open-irreducibility}; normalization by
the inverse cyclotomic character does not change it.
The proposition therefore verifies the infinite-extension
irreducibility condition required by
\href{https://arxiv.org/pdf/2207.04925}{Thorne, Theorem 6.1},
once the specified regular algebraic polarized cuspidal
automorphic realization has been constructed.
That theorem concerns the adjoint descended to the maximal
totally real subfield.
Its application still requires the actual tangent classes to
satisfy the Bloch--Kato local conditions.
The proposition supplies the irreducibility condition, rather
than the automorphic realization or those local conditions.

\subsection{A fixed finite catalogue of imaginary quadratic fields}
\label{sec:sr-finite-CM-catalogue}
The rank-four automorphy predicate used by the source provides
automorphy over imaginary quadratic fields where the restriction
is irreducible. It does not itself assert automorphy over
\(\mathbb Q\). The following elementary argument selects such a
field from a fixed finite list.

\begin{srlemma}\label{lem:sr-quadratic-self-twists}
Let \(G\) be a group and \(V\) an irreducible representation of
dimension \(n\) over an algebraically closed field of
characteristic zero.
There are at most \(n^2\) characters \(\eta:G\to L^\times\)
with \(V\simeq V\otimes\eta\), including the trivial character.
If \(H\) is normal of index two with quadratic character
\(\eta_H\), then \(V|_H\) is reducible exactly when
\(V\simeq V\otimes\eta_H\).
An irreducible representation of \(H\) has at most two isomorphism
classes of extensions to \(G\).
\end{srlemma}
\begin{proof}
A self-twist gives a nonzero intertwiner in \(\operatorname{End}_L(V)\);
it is invertible by irreducibility.
It spans an eigenline for conjugation by \(G\), with character
\(\eta\) or its inverse according to the intertwiner convention.
Different characters give linearly independent eigenlines.
To prove this directly, take a linear relation of minimal
nonzero length among chosen eigenvectors. Conjugate by an element
separating two of its characters and subtract the appropriate
scalar multiple of the original relation. This gives a nonzero
relation of smaller length, a contradiction.
Their number is at most \(\dim\operatorname{End}_L(V)=n^2\).

The restriction to a normal subgroup is semisimple by the
simple-submodule orbit argument in the preceding proof.
For an index-two subgroup its isotypic components form one orbit
of size one or two.
If there is one type with multiplicity greater than one,
the \(T\otimes B\) argument for a lift of the nontrivial coset
gives a proper \(G\)-invariant subspace from an eigenline of \(B\).
This contradicts irreducibility.
Thus reducibility of \(V|_H\) forces two isotypic components.
The nontrivial coset swaps them.
The operator equal to \(1\) on the first and \(-1\) on the second
commutes with \(H\) and anticommutes with that coset.
It intertwines \(V\) with \(V\otimes\eta_H\).

Conversely, let \(T\) be such an intertwiner.
Its square commutes with \(G\), hence is a nonzero scalar.
Rescale \(T\) so that \(T^2=1\).
The operator \(T\) cannot be scalar, since it anticommutes with an
invertible operator in the nontrivial coset and the characteristic
is not two.
Both its \(1\) and \(-1\) eigenspaces are therefore nonzero,
and both are \(H\)-invariant. This proves the converse.

Finally identify the restrictions of two extensions to an
irreducible \(H\)-representation.
For a fixed element \(g\) outside \(H\), their two operators
implement the same conjugation action on \(H\).
Their ratio commutes with \(H\), and so is scalar by Schur's lemma.
Their squares agree because \(g^2\in H\), so that scalar is \(1\)
or \(-1\). These are precisely the two twists by the quotient
characters. This proves the last assertion.
\end{proof}

Choose \(n^2+1\) distinct rational primes \(p_i\equiv3\pmod4\)
and put \(K_i=\mathbb Q(\sqrt{-p_i})\).
These choices are made before the conductor, twist and coefficient
prime. Only finitely many primes are needed; their existence follows,
for example, from Dirichlet's theorem.
For any continuous irreducible rank-\(n\) representation of
\(G_{\mathbb Q}\), at most \(n^2-1\) of these restrictions are
reducible, by the lemma.
For an odd coefficient prime \(\ell\), at most one \(K_i\) is
ramified at \(\ell\), since its discriminant is \(-p_i\).
Thus at least one fixed \(K_i\) has irreducible restriction and
is unramified at \(\ell\).
For rank four a list of seventeen fields suffices.

There is a corresponding finite-family counting statement.
Assign each member of a finite set of pairwise nonisomorphic
irreducible \(G_{\mathbb Q}\)-representations to one field from
the fixed list where it restricts irreducibly.
If the restrictions assigned to \(K_i\) have at most \(B_i\)
isomorphism classes, the last assertion of the lemma gives
\[
                         |\mathscr R|\leq2\sum_i B_i.
\]
In particular constants in established counts over the fixed
fields can be maximized before choosing the varying host.
This is a transfer of an already verified restriction count.
Applying an automorphic count also requires its precise
weight and level conditions for these restrictions.

\subsection{A rank-four cyclotomic irreducibility criterion}
\label{sec:sr-rank-four-cyclotomic}

The following criterion supplies the cyclotomic irreducibility part of
the rank-four rigidity argument. It uses the residual quadratic
self-twist exclusion, rather than an irreducibility assertion on all
open subgroups.

\begin{srlemma}\label{lem:sr-rank-four-no-self-twist}
Let \(\ell\) be odd, let \(\Gamma\) be profinite, and let
\(\epsilon:\Gamma\to\mathbb Z_\ell^\times\) be continuous.
Let \(V\) be a continuous absolutely irreducible representation of
dimension four over a finite extension of \(\mathbb Q_\ell\).
Suppose a stable lattice has residual semisimplification
\(\chi_1\oplus\cdots\oplus\chi_4\), over a finite splitting residue
field \(k\). Write \(\bar\epsilon\) for the residual cyclotomic
character. Assume
\[
 \left.
 \begin{gathered}
 \eta:\Gamma\to k^\times,\quad \eta^2=1,\quad
            \ker\bar\epsilon\subseteq\ker\eta,\\
 \{\eta\chi_1,\ldots,\eta\chi_4\}
                     =\{\chi_1,\ldots,\chi_4\}
                   \quad\hbox{as multisets}
 \end{gathered}\right\}
                  \quad\Longrightarrow\quad\eta=1.
 \tag{C}
\]
Then \(V|_{\ker\epsilon}\) is absolutely irreducible.
\end{srlemma}
\begin{proof}
Extend scalars to an algebraic closure \(L\) of the coefficient
field, and put \(H=\ker\epsilon\).
The quotient \(\Gamma/H\) is a closed subgroup of
\(\mathbb Z_\ell^\times\), hence is procyclic for odd \(\ell\).
Choose \(g\in\Gamma\) lifting a topological generator.

The restriction to \(H\) is semisimple.
Indeed choose an \(H\)-submodule \(W\ne0\) of smallest dimension.
It is simple, and its \(\Gamma\)-translates are simple \(H\)-submodules
because \(H\) is normal. Their sum is nonzero and
\(\Gamma\)-invariant, so is all of \(V_L\).
A finite subcollection spans \(V_L\).
A finite sum of simple modules is semisimple: when a further simple
module is added, its intersection with the preceding sum is either
zero or that whole simple module.

Let \(V_0,\ldots,V_{d-1}\) be the isotypic components.
The action of \(\Gamma\) on them is transitive; otherwise the sum
over one orbit is a proper \(\Gamma\)-submodule.
Each stabilizer is closed, by the vanishing of a block of matrix
entries in a basis adapted to that component.
It has finite index and is therefore open.
The permutation action consequently factors continuously through
the procyclic quotient \(\Gamma/H\).
Thus \(g\) cyclically permutes the components.
They have equal dimension, so \(d\mid4\).

If \(d>1\), then \(d=2\) or \(4\).
Let \(S\) be the stabilizer of \(V_0\). Its image in
\(\Gamma/H\) has cyclic quotient of order \(d\).
The character of this quotient taking \(g\) to \(-1\) defines a
continuous nontrivial quadratic character
\(\eta_0:\Gamma\to\{1,-1\}\).
The operator \(T\) which is \((-1)^i\) on \(V_i=g^iV_0\)
commutes with \(S\) and anticommutes with \(g\).
Therefore \(V_L\simeq V_L\otimes\eta_0\).

Any continuous quadratic character of a closed subgroup of
\(\mathbb Z_\ell^\times\) kills its intersection with
\(1+\ell\mathbb Z_\ell\): that intersection is pro-\(\ell\),
whereas its image has order at most two.
Hence \(\eta_0\) factors through \(\bar\epsilon\).
Reduce its values to \(k\), obtaining \(\bar\eta_0\).
The self-twist gives equality of characteristic polynomials on
every element; the residual semisimplifications therefore have
the same Jordan--Hölder constituents, by the characteristic-polynomial
form of Brauer--Nesbitt. Consequently \(\bar\eta_0\) permutes the
multiset of the four residual characters.
Condition (C) gives \(\bar\eta_0=1\).
Reduction is injective on \(\{1,-1\}\) because \(\ell\ne2\),
contradicting the choice of \(\eta_0\).
Thus \(d=1\).

Write the single isotypic restriction as
\(V_L=W\otimes_L M\), with \(W\) simple for \(H\).
Conjugation by \(g\) preserves the isomorphism class of \(W\).
Choose an intertwiner \(A\) for this conjugation.
After multiplying the action of \(g\) by \(A^{-1}\) on \(W\),
the resulting operator commutes with \(H\).
Schur's lemma gives
\(\operatorname{End}_H(V_L)=1_W\otimes\operatorname{End}_L(M)\);
hence \(g\) acts as \(A\otimes B\), for some invertible \(B\) on \(M\).
An eigenline \(M_1\) of \(B\) makes \(W\otimes M_1\) invariant
under \(H\) and \(g\).

The subgroup generated by \(H,g\) is dense in \(\Gamma\).
The stabilizer of the linear subspace \(W\otimes M_1\) is closed,
again by block-matrix equations, so this subspace is
\(\Gamma\)-invariant. Absolute irreducibility forces
\(W\otimes M_1=V_L\). Thus \(\dim M=1\), proving the result.
\end{proof}

The criterion (C) is exactly the residual exclusion in
\(\mathrm{TNiceFour.nice0four\_alt}\), called
\(\mathrm{P1.condC}\) at the linked point.
The latter is proved in
\(\mathrm{WeakShimuraTheory.ProjectiveLineFlags.CondC}\).
The four characters in the actual linked packet are pairwise distinct
by the tower's separation condition at its distinguished element.
Thus the source's inclusion condition on the four characters is
equivalent here to permutation of the multiset.
Its transport to the linked point involves no estimate:
the level quotient map is surjective, its residual cyclotomic
character is the usual one by \(\mathrm{P1.levEps}\), and its
four residual characters have a common scalar-character factor.
That common factor cancels from a proposed permutation.
A quadratic character takes only the values \(1,-1\), so it descends
through an injective residue-field extension.
Pulling it back, cancelling the common factor, and using the
tower's residual self-twist exclusion gives (C).
This explains exactly which residual property of the constructed
tower is being used. It does not construct that tower.

Applied to the linked rank-four point, the lemma proves
irreducibility over the full cyclotomic kernel.
Scalar twists do not change invariant subspaces.
Surjectivity of the level quotient identifies its cyclotomic
kernel with the image of \(G_{L(\mu_{\ell^\infty})}\).
Thus the same assertion holds for the underlying restricted
Galois representation. The additional condition called
\(\mathrm{P1.condD}\), and the source's regular-coset property
\(\mathrm{N3}\), are not needed for this particular deduction.
They have separate uses in the source's patching argument.

\subsection{The local condition for an ordinary first-order lift}
\label{sec:sr-rank-four-local-tangent}

We give the local reason for the Bloch--Kato condition; ordinary
triangularity by itself would not suffice at an exceptional adjacent
character ratio.

\begin{srlemma}\label{lem:sr-positive-ordinary-extension}
Let \(F/\mathbb Q_\ell\) be finite and \(E/\mathbb Q_\ell\) a finite
coefficient extension. Let
\(\widetilde r:G_F\to\operatorname{GL}_n(E[\tau]/\tau^2)\)
be continuous and upper triangular.
Its diagonal characters are \(\widetilde\kappa_i\), with reductions
\(\kappa_i\).
Suppose there are strictly decreasing integers
\(w_1>\cdots>w_n\) such that on inertia
\[
                \widetilde\kappa_i=\epsilon^{w_i}
\]
as \(E[\tau]/\tau^2\)-valued characters.
Assume also that
\[
               \kappa_i\kappa_j^{-1}\ne\epsilon
                  \quad\hbox{on }G_F\quad (i<j).
 \tag{N}
\]
Then \(\widetilde r\), regarded as a \(2n\)-dimensional
\(E\)-representation, is crystalline.
Its reduction \(r\) is crystalline, and its first-order class
belongs to \(H^1_f(F,\operatorname{ad}r)\).
\end{srlemma}
\begin{proof}
The \(i\)-th diagonal block, regarded over \(E\), is a
two-dimensional unramified self-extension of
\(\kappa_i=\epsilon^{w_i}\alpha_i\), with \(\alpha_i\) unramified.
It is crystalline; unramified representations are crystalline
and tensoring by an integral cyclotomic power preserves that property.

Build the upper triangular representation one diagonal block at a time.
Suppose the preceding blocks form a crystalline representation \(U\),
and the next block is \(W\), of weight \(w_j\).
The representation \(T=\operatorname{Hom}_E(W,U)\) is crystalline.
Its semisimplification consists of the characters
\(\kappa_i\kappa_j^{-1}\), for \(i<j\), with their multiplicities.
Their cyclotomic exponents \(w_i-w_j\) are positive.
Thus \(H^0(F,T)=0\) and
\[
      \dim_E D_{\rm dR}(T)/\mathrm{Fil}^0
                          =[F:\mathbb Q_\ell]\dim_E T.
\]
Moreover \(H^0(F,T^*(1))=0\): if \(w_i-w_j>1\), its corresponding
character has nonzero cyclotomic exponent; if \(w_i-w_j=1\),
a fixed vector would contradict (N).
A representation with a finite filtration and no trivial
Jordan--Hölder constituent has no fixed vector, so this check on
the constituents suffices.

Local Tate duality gives \(H^2(F,T)=0\).
The local Euler characteristic and the crystalline fundamental
exact sequence give respectively
\[
 \dim_E H^1(F,T)=[F:\mathbb Q_\ell]\dim_E T,\qquad
 \dim_E H^1_f(F,T)
       =\dim_E D_{\rm dR}(T)/\mathrm{Fil}^0+\dim_E H^0(F,T).
\]
The dimensions agree. Since \(H^1_f(F,T)\subseteq H^1(F,T)\),
every extension class is crystalline.
Therefore the extension of \(W\) by \(U\) in \(\widetilde r\)
is crystalline. Induction proves the assertion.
Finally the sequence
\(0\to r\to\widetilde r\to r\to0\), regarded over \(E\),
identifies its class with the defining crystalline local
condition in \(H^1(F,\operatorname{ad}r)\).
\end{proof}

Here is the concrete check for the rank-four packet.
Write the ordinary cyclotomic exponents as
\[
                         (a+1,a,1,0),\qquad a\geq2.
\]
Choose the source's Frobenius lift \(\sigma\) at \(\ell\) with
\(\epsilon(\sigma)=1\).
On the ordered residual flag its eigencharacters at \(\sigma\)
have values
\[
                              uv,\quad u,\quad v,\quad1,
\]
where \(u=\bar\psi(\sigma)\) and \(v=\bar\xi(\sigma)\).
The good-data condition is
\[
                              vu^m\ne1\qquad(m\in\mathbb Z).
 \tag{D}
\]
For weight difference one, the first and second positions, and the
third and fourth positions, have residual ratio \(v\ne1\).
When \(a=2\), the middle positions also have weight difference one,
and their residual ratio is \(u/v\ne1\), by (D) with \(m=-1\).
These are all the pairs with weight difference one.

The remaining finite inertial characters can be killed over a
finite local extension whose Galois subgroup still contains \(\sigma\).
To see this, collect those characters into one finite-image character
\(\chi\). Divide \(\chi\) by the unramified character taking
\(\sigma\) to \(\chi(\sigma)\), and take its kernel.
This is open and normal, contains \(\sigma\), and kills \(\chi\)
on inertia. Over this extension every ordinary diagonal has the
indicated pure cyclotomic restriction on inertia.
A ratio with weight difference other than one cannot become
\(\epsilon\) on an open group, since the cyclotomic image on inertia
is infinite. For difference one, the residual check at the retained
\(\sigma\) excludes that equality.
Thus Lemma~\ref{lem:sr-positive-ordinary-extension} applies.

Crystalline local vanishing descends from this finite extension.
For the period-module cohomology defining \(H^1_f\), restriction
is injective after replacing the extension by a finite normal one:
the inflation--restriction kernel is a degree-one cohomology group
of a finite group acting on a characteristic-zero vector space,
and averaging kills it.
The argument therefore proves the original Bloch--Kato condition.

Away from \(\ell\), the source's local conditions make the lift
trivial on an open subgroup of inertia.
The tangent cocycle is then zero on that subgroup.
An explicit finite-index averaging proves it is an inertia
coboundary: after a local coboundary is subtracted, let the cocycle
be \(z\), zero on \(J\leq I\) of finite index \(m\).
It factors as a function on left cosets, and if
\(S=\sum_{gJ\in I/J}z(g)\), the cocycle identity gives
\(S=mz(h)+hS\). Hence \(z(h)=(h-1)(-S/m)\).

For the formal point problem, the ordered flags and their residual
labels are the ones specified by the local deformation conditions.
One first descends the point and its dual-number lift to a finite
coefficient field, as in the source's
\(\mathrm{TbTriDiag.tb\_desc}\).
The preceding proof then supplies the local content of
\(\mathrm{TbTriDiag.tanBK}\), including the adjacent-ratio check.
It does not infer the Bloch--Kato condition solely from ordinary
triangularity.

\subsection{The correct polarized adjoint and the symplectic projection}
\label{sec:sr-rank-four-adjoint-projection}

\begin{srproposition}\label{prop:sr-polarized-symplectic-rigidity}
Let \(r:G_{\mathbb Q}\to\operatorname{GSp}_4(E)\) be continuous,
with fixed similitude character \(\theta\), and let \(K\) be an
imaginary quadratic field with character \(\eta\).
Fix a complex conjugation \(c\) with \(\theta(c)=-1\).
Let \(J\) be the alternating matrix preserved up to \(\theta\), and set
\[
 X^\#=J^{-1}X^{\mathsf T}J,\qquad
 P(X)=\tfrac12(X-X^\#),\qquad
 A(g)X=r(g)Xr(g)^{-1}.
\]
Define the representation on \(M_4(E)\) by
\[
 R_\eta(g)X=A(g)P(X)+\eta(g)A(g)(X-P(X)).
 \tag{R}
\]
It is the polarized adjoint descended from \(r|_{G_K}\), with
polarizing matrix \(B=Jr(c)^{-1}\):
\[
 R_\eta(h)X=A(h)X\quad(h\in G_K),\qquad
 R_\eta(c)X=-B^{-1}X^{\mathsf T}B.
\]
Suppose its Bloch--Kato Selmer group over \(\mathbb Q\) vanishes.
Then every fixed-similitude first-order lift of \(r\) whose tangent
class satisfies the local conditions proved above is conjugate
to the constant lift by \(1+\tau Y\), with \(Y\in\mathfrak{sp}_4(E)\).
\end{srproposition}
\begin{proof}
The symplectic adjoint is an involution and commutes with
conjugation by a symplectic similitude. Thus \(P^2=P\),
its image is \(\mathfrak{sp}_4(E)\), and
\(P A(g)=A(g)P\).
The two invariant summands \(\operatorname{im}P\) and \(\ker P\)
show directly that (R) is a representation.
On \(G_K\), \(\eta=1\), giving the ordinary adjoint action.
At \(c\), \(\eta=-1\), so
\[
 R_\eta(c)X=A(c)(2P(X)-X)=-A(c)X^\#
                                      =-B^{-1}X^{\mathsf T}B.
\]
For \(h'=chc^{-1}\), the similitude identity gives
\[
                  r(h)^{\mathsf T}B\,r(h')=\theta(h)B,
\]
which identifies the asserted polarization and its descent.
Since \(r(c)^2=1\) and \(\theta(c)=-1\), the same identity gives
\(B^{\mathsf T}=B\), the required odd polarized descent in this
rank-four setting.

Write a first-order lift as
\(\widetilde r(g)=r(g)+\tau\,\dot r(g)\), and put
\(z(g)=\dot r(g)r(g)^{-1}\).
The homomorphism identity says
\(z(gh)=z(g)+A(g)z(h)\).
Differentiating the fixed-similitude identity
\(\widetilde r(g)^{\mathsf T}J\widetilde r(g)=\theta(g)J\)
gives \(z(g)^\#=-z(g)\). Hence \(Pz=z\).

The map \(P\) intertwines \(A\) with \(R_\eta\):
\(P A(g)=R_\eta(g)P\).
It carries inertia-coboundary witnesses to inertia-coboundary
witnesses, and its tensor product with the identity on
\(B_{\rm cris}\) does the same for the crystalline witnesses.
Thus \(z\) is a Bloch--Kato cocycle for \(R_\eta\).
By the assumed vanishing it is a coboundary,
\(z(g)=R_\eta(g)X-X\).
Projecting once more gives
\[
                    z(g)=A(g)P(X)-P(X).
\]
Set \(Y=-P(X)\). Then \(Y\in\mathfrak{sp}_4(E)\) and
\[
                    \dot r(g)=Yr(g)-r(g)Y.
\]
This is exactly conjugation by \(1+\tau Y\).
\end{proof}

The applicable published vanishing statement is
\href{https://arxiv.org/pdf/2207.04925}{Thorne, Theorem 6.1}:
for an actual regular algebraic polarized cuspidal realization
over a CM field, cyclotomic-kernel irreducibility implies
vanishing of its polarized adjoint descended to the maximal
totally real subfield.
This is not a statement about the full adjoint Selmer group
over the CM field.

For a realization over a larger CM field \(L\supset K\), the
same application is valid after the polarization is matched.
Since \(K\cap L^+=\mathbb Q\), one has \(L=KL^+\).
Consequently \(R_\eta|_{G_{L^+}}\) is precisely the descended
adjoint of \(r|_{G_L}\) with the displayed polarization.
An algebraic scalar twist does not change its adjoint on \(G_L\);
its polarizing form must still be identified, up to scalar,
with the form defining this descent.
This is the purpose of the scalar norm identity in
\(\mathrm{AdDown}\), not a consequence of adjoint
twist-invariance alone.
Once this matching and the actual polarized cuspidal realization
are supplied, the vanishing over \(L^+\) implies the vanishing
over \(\mathbb Q\): restriction preserves the local conditions
and is injective in degree one by finite-index averaging in
characteristic zero.

The rank-four automorphy predicate \(\mathrm{IsAut4Alt}\) is
\(\mathrm{IsAutQA}\). Its exact definition provides a regular
algebraic cuspidal realization over every imaginary quadratic
field where the restriction is irreducible, together with
\(\mathrm{WitA}\). It does not assert a cuspidal
\(\operatorname{GL}_4/\mathbb Q\) realization.
The field selection, local level compatibility and polarization
matching in the preceding paragraph are therefore essential.

\paragraph{What these arguments close.}
They prove the residual criterion for cyclotomic irreducibility,
the ordinary tangent's Bloch--Kato local condition with its
nonexceptionality check, and the symplectic projection from
the correctly descended polarized adjoint.
They do not prove the quantitative exceptional estimate without
the actual rank-four finite/no-block/automorphic-point producer.
That source producer is
\(\mathrm{Sp4In0.sp4In0\_alt}\), through
\(\mathrm{CuFinal.main4\_alt}\) and
\(\mathrm{ModularityRankFourOfPromodularity}\).
Its Case II input is
\(\mathrm{Gsp4PromodularityCaseTwo.proMod\_of\_pmEx\_of\_ctx2\_of\_kPair\_of\_star}\),
which constructs the deep residue chain and propagates the modular
support from an automorphic seed.
For the present principal packet, no Case III count is needed.
The Gaussian twist-one branch still needs the marked universal-ring
and type-I automorphy statement \(\mathrm{ThmS6Core\_s}\), supplied
in the formal dependency chain by \(\mathrm{DnEndsTop.ends0\_alt}\).
All paths in this paragraph are in the frozen tag
\texttt{alpoge-source-7cbf7a0c}.
